\begin{textsample}{NEG Dim 3 – summary\_guided\_gpt – Score -3.00 – summary\_guided\_gpt/t2707\_gpt.txt}  \label{ex:f3_neg_024}
I ’ve always been really comfortable being alone, and for a long time I thought that meant there was something wrong with me.
Everyone around me seems to \textbf{crave} company, to \textbf{feel} energized by hanging out, texting, calling, making plans.
Meanwhile, when my friends reach out – people I genuinely like and care about – my first reaction is anxiety.
I start \textbf{worrying} about how to “ socialize correctly, ” what to say, how long I ’m supposed to hang out, whether I ’ll come off as weird or distant.
Even if I do meet up, I usually end up feeling overwhelmed pretty quickly and just want to go home and be by myself again.
It ’s confusing because it ’s not that I hate people or that I don't want friends.
I do.
I just \textbf{feel} this huge sense of relief when plans get cancelled or when a weekend passes and no one expects anything from me.
The COVID lockdown really highlighted this.
While a lot of people were struggling with isolation, I \textbf{felt} … relieved.
Staying home, not having to contact anyone, not \textbf{feeling} pressured to “ keep up ” socially – it was like I could finally breathe.
That made me wonder if this is something more extreme than just being introverted.
Finding this community has been a weird \textbf{mix} of comforting and emotional.
Reading posts from people who \textbf{feel} the same way makes me \textbf{feel} understood in a way I honestly didn't expect.
For once I don't \textbf{feel} like a broken friend or a bad person for needing so much space.

% matched lemmas: crave, feel, mix, worry
\end{textsample}
